\documentclass{ctexart}
\usepackage{amsmath,amssymb}  % 数学公式基础包
\usepackage{listings}         % 代码排版包
\usepackage{xcolor}           % 代码高亮颜色

% 代码样式配置（适配大一学生阅读）
\lstset{
    language=C++,
    basicstyle=\small\ttfamily,
    keywordstyle=\color{blue},
    commentstyle=\color{green!60!black},
    stringstyle=\color{red},
    numbers=left,
    numberstyle=\tiny\color{gray},
    frame=single,
    breaklines=true,
    showspaces=false,
    showstringspaces=false
}

% 自定义数学符号（简化书写）
\newcommand{\ii}{\mathrm{i}}       % 虚数单位
\newcommand{\e}{\mathrm{e}}        % 自然常数
\newcommand{\F}{\mathcal{F}}       % 傅里叶变换算子
\newcommand{\om}[1]{\omega_{#1}}   % 单位根简写

\title{快速傅里叶变换（FFT）—— 线性代数视角（精简版）}
\author{适配大一线性代数项目（原作者：H-J-Granger 等）}
\date{\today}

\begin{document}
\maketitle

\section{前置知识}
复数的基本运算、矩阵与线性变换、多项式的系数/点值表示。
精要说明：下面不提供长段程序代码，仅保留对算法要点的清晰阐释与极短伪代码，便于理解原理后自行查阅或实现。

\subsection{算法要点}
- 变换目标：将多项式的系数表示转换为在 $N$ 个单位根处的点值表示（DFT），反之为 IDFT。\
- 单位根性质：$\omega_N^{k+N/2}=-\omega_N^k$，使得在合并阶段存在配对对称性。\
- 分治策略：把长度为 $N$ 的序列分为偶、奇两部分，分别对长度 $N/2$ 的子序列递归做 DFT，再按公式合并。\
- 复杂度：递推式 $T(N)=2T(N/2)+O(N)$，得 $T(N)=O(N\log N)$。\
- 逆变换：用相反的指数符号并在最终结果除以 $N$ 完成归一化。

\subsection{数学合并公式（核心合并步）}
设 $G(k)$ 与 $H(k)$ 分别为偶、奇子序列对 $\omega_{N}^{2k}$ 的 DFT 结果，则：
\[
F(k)=G(k)+\omega_N^{k}H(k),\qquad F(k+N/2)=G(k)-\omega_N^{k}H(k),
\]
其中 $\omega_N=\mathrm{e}^{-2\pi\mathrm{i}/N}$（DFT 情形），逆变换取相反指数并整体除以 $N$。

\section{目标与背景}
快速傅里叶变换（FFT）的目标是高效计算长度 $N$ 的离散傅里叶变换（DFT）：
\[
X_k=\sum_{n=0}^{N-1} x_n\,\omega_N^{kn},\qquad\omega_N=\mathrm{e}^{-2\pi\mathrm{i}/N},\;k=0,\dots,N-1.
\]
直接计算需要 $O(N^2)$ 次复乘加，FFT 的目标是将复杂度降低到 $O(N\log N)$。

\section{单位根与对称性}
核心代数事实是单位复根的周期性与对称性：$\omega_N^N=1$，并且当 $N$ 为偶数时有
\[
\omega_N^{k+N/2}=-\omega_N^k.
\]
该等式导致在某些评价点处的函数值成对出现符号关系，从而可以重用子问题的计算结果。

\section{分治思想的代数推导}
将序列 $x=(x_0,x_1,\dots,x_{N-1})$ 按偶/奇下标拆为两部分：偶项 $a_m=x_{2m}$ 与奇项 $b_m=x_{2m+1}$。构造多项式
\[
f(z)=\sum_{n=0}^{N-1} x_n z^n = G(z^2)+zH(z^2),
\]
其中 $G(y)=\sum_m a_m y^m$，$H(y)=\sum_m b_m y^m$。在点 $z=\omega_N^k$ 处评价得：
\[
f(\omega_N^k)=G(\omega_N^{2k})+\omega_N^k H(\omega_N^{2k}).
\]
注意到 $\omega_N^{2}=\omega_{N/2}$，因此 $G(\omega_N^{2k})$ 和 $H(\omega_N^{2k})$ 都是对长度 $N/2$ 的 DFT 在索引 $k$ 处的值。

此外，对于配对的点 $\omega_N^{k+N/2}$，有
\[
f(\omega_N^{k+N/2})=G(\omega_{N/2}^k)+\omega_N^{k+N/2}H(\omega_{N/2}^k)=G(k)-\omega_N^{k}H(k).
\]
将以上两式并列写出合并公式：对每个 $k=0,\dots,N/2-1$，
\[
F(k)=G(k)+\omega_N^{k}H(k),\qquad F(k+N/2)=G(k)-\omega_N^{k}H(k).
\]
因此，计算长度 $N$ 的 DFT 等价于计算两个长度 $N/2$ 的 DFT，再进行 $O(N)$ 的合并。

\section{复杂度分析}
设 $T(N)$ 为长度 $N$ 的 DFT 运算量（按复乘加计），由分治关系：
\[
T(N)=2T(N/2)+cN.
\]
按主定理得 $T(N)=O(N\log N)$，这就是 FFT 将计算量从 $O(N^2)$ 降为 $O(N\log N)$ 的数学来源。

\section{矩阵与代数分解视角}
DFT 可视为乘以 Vandermonde 型矩阵 $F_N$：$X=F_N x$，其中 $(F_N)_{k,n}=\omega_N^{kn}$。
当 $N$ 可因式分解（尤其 $N=2^m$）时，$F_N$ 可分解为块矩阵与对角旋转因子的乘积，代数上将密集矩阵分解为若干结构化稀疏块，从而减少实际乘法次数。

\section{卷积定理与多项式乘法}
卷积定理指出离散卷积在频域对应逐点相乘：
\[
\mathrm{DFT}(a*b)=\mathrm{DFT}(a)\odot\mathrm{DFT}(b).
\]
因此多项式系数的卷积（乘法）可通过 DFT → 点值相乘 → IDFT 完成，整体复杂度为 $O(N\log N)$。

\section{逆变换与归一化}
逆离散傅里叶变换（IDFT）给出系数的恢复公式：
\[
x_n=\frac{1}{N}\sum_{k=0}^{N-1} X_k\,\omega_N^{-kn}.
\]
在分治合并中使用相反的指数并在最终结果除以 $N$ 即完成逆变换的归一化步骤。

\section{数值与实现要点（纯描述）}
- 常见优化：将输入补零到 $2^m$ 长度以便递归分解；在工程实现中使用迭代原地算法与位反转重排以节省内存。\\ 
- 精度问题：浮点运算会带来舍入误差，通常通过使用双精度（double）与合理的缩放/归一化策略减小误差。\\ 
- 对称性利用：实值输入可利用共轭对称性减少计算；旋转因子可预计算并复用以节省成本。

\section{4 点手算示例（直观）}
设 $N=4$，序列 $x=(x_0,x_1,x_2,x_3)$，单位根 $\omega_4=\mathrm{e}^{-2\pi\ii/4}=-\mathrm{i}$（取惯例符号）。拆分得到：
- 偶项多项式 $G$（长度 2）：$A_0=x_0+x_2,\;A_1=x_0-x_2$.\\ 
- 奇项多项式 $H$（长度 2）：$B_0=x_1+x_3,\;B_1=x_1-x_3$。

合并得到：
\begin{align*}
F(0)&=A_0+\omega_4^{0}B_0=(x_0+x_2)+(x_1+x_3)=x_0+x_1+x_2+x_3,\\
F(2)&=A_0+\omega_4^{2}B_0=(x_0+x_2)-(x_1+x_3)=x_0-x_1+x_2-x_3,\\
F(1)&=A_1+\omega_4^{1}B_1=(x_0-x_2)+\omega_4(x_1-x_3),\\
F(3)&=A_1+\omega_4^{3}B_1=(x_0-x_2)+\omega_4^{3}(x_1-x_3).
\end{align*}
此例说明如何通过两次长度 2 的小变换与少量合并运算得到长度 4 的 DFT 结果。

\section{小结}
FFT 的核心在于利用单位根的周期与对称性，将密集的 DFT 变换分解为可重用的子变换与简单的合并步骤，从而在代数层面与算法层面同时实现效率提升。
\end{document}